% 中文
\cabstract{
根据实验内容 \cite{lab03}, 在此次实验需要实现一个基于 Java 的简单中缀表达式到后缀表达式的转换程序（Postfix），此报告参照评分标准\cite{labeval03}，完成这份实验报告撰写并附上相应的实验报告说明。

在本次 Postfix 项目编译分析器的重修与升级实验中，
对静态成员变量进行了深入分析，发现其在多实例或多线程环境下会引发状态污染和解析异常，因此将静态成员变量修改为实例变量以提升程序的健壮性和适用范围；
并在不损耗 $O(n)$ 时间复杂度的前提下，用循环迭代替代了递归，解决栈帧溢出的问题。
从原本容误度极低的简单错误判断，进一步升级为了有能力探测所有文法语缀差错点的错误分析工具；
并且搭建 JUnit 黑盒测试层配合完整的自动化 Pipeline 等实战工程化环节，
此次实验拔高了我统筹编写高泛性、深稳健型实用底层程序的代码能力。
}

\ckeywords{编译原理；中缀转后缀；尾递归消除；错误恢复；Junit单元测试}